% TOP_PROBLEM: {"id": 30007120, "title": "Logarithmic Sobolev inequality from subgaussian marginals", "category_id": 6, "category": {"id": 6, "name": "geometry", "display_name": "Geometry", "description": "Euclidean and non-Euclidean geometry, geometric structures.", "slug": "geometry", "order_index": 6, "created_at": "2026-07-31T15:26:25.671Z"}, "status": "open"}
% FIRST_POSTED: 2026-10-06
% ATTEMPT_STATUS: solved
% OPENAI_CLAIM: {"schema_version":1,"source_repo":"https://github.com/openai/math","source_commit":"adc7f1241b42e322a6451854ab7e4b4c146bf78a","release_date":"2026-10-06","match":"full","resolution":"proved","independently_reviewed":false,"second_pass":"agreed","families":[{"family":"093","title":"Dimension-free logarithmic Sobolev inequality for subgaussian log-concave measures","subject":"Convex and metric geometry","manuscripts":[{"title":"A dimension-free logarithmic Sobolev inequality for subgaussian log-concave measures","pdf_path":"preprints/A-dimension-free-logarithmic-Sobolev-inequality-for-subgaussian-log-concave-measures-September-23-2026/paper.pdf","readme_path":"preprints/A-dimension-free-logarithmic-Sobolev-inequality-for-subgaussian-log-concave-measures-September-23-2026/README.md","date":"2026-09-23","theorem_ref":"thm:main (sections/01-introduction.tex:26)","statement_sources":[{"path":"preprints/A-dimension-free-logarithmic-Sobolev-inequality-for-subgaussian-log-concave-measures-September-23-2026/build/sections/01-introduction.tex","line":26,"label":"thm:main"}]}],"lean":null,"reasoning_trace":null}]}
\documentclass[11pt]{article}
\usepackage[margin=25mm]{geometry}
\usepackage{fontspec}
\setmainfont{Latin Modern Roman}
\usepackage{amsmath,amssymb,mathtools}
\usepackage{unicode-math}
\setmathfont{Latin Modern Math}
\usepackage{enumitem}
\usepackage{xurl,hyperref,bookmark}
\hypersetup{colorlinks=true,urlcolor=blue}
\setcounter{secnumdepth}{0}
\setlength{\parindent}{0pt}
\setlength{\parskip}{5pt}
\setlength{\emergencystretch}{3em}
\begin{document}
\section*{Logarithmic Sobolev inequality from subgaussian marginals}

\subsection{OpenAI's claim}
OpenAI presents the following claims relevant to this question, with manuscript numbers following paper order in Sources. In manuscript 1 (family 093; 2026-09-23), the asserted result is that a centered log-concave density with linear subgaussian parameter $a$ satisfies a logarithmic Sobolev inequality with constant bounded by an absolute multiple of $a^2$.

\subsection{Mathematical statement of the claimed result}
\paragraph{Manuscript 1.}
Take any $n\ge1$, a centered log-concave probability measure $\mu$ on $\mathbb R^n$ with Lebesgue density, and $a>0$ such that $X\sim\mu$ obeys
\[
 \sup_{|\theta|=1}\mathbb E\exp(\langle X,\theta\rangle^2/a^2)\le2.
\]
The claimed absolute, dimension-independent constant $C$ gives, for all $f\in C_c^{\infty}(\mathbb R^n)$,
\[
 \operatorname{Ent}_{\mu}(f^2)\le Ca^2\int|\nabla f|^2\,d\mu,
 \qquad \operatorname{Ent}_{\mu}(g)=\int g\log g\,d\mu-
 (\int g\,d\mu)\log(\int g\,d\mu).
\]
Only linear-function subgaussianity is assumed, but a logarithmic Sobolev conclusion is claimed.

\subsection{Outline of the argument}
For manuscript 1 (family 093; 2026-09-23), gaussian regularization and entropy dissipation reduce the question to critical posterior laws. Efficient observation probes and nearly-first-eigenfunction hierarchies control the prediction curve, while curvature and operator-tail bounds rule out a dimension-dependent obstruction.

\subsection{Scope relative to this problem}
The result supplies the conclusion in this problem with its displayed hypotheses and quantifiers. The broader statements and applications in the manuscript are kept separate from the particular question above. The supporting result is manuscript 1 (family 093; 2026-09-23).

\subsection{Formal verification}
No Lean formalization is listed at this snapshot for manuscript 1 (family 093; 2026-09-23).

This site has not executed the Lean code. Described statement coverage is not an independent proof check.

\subsection{Status caveat}
The manuscripts were produced by an internal OpenAI model. OpenAI cautions that results without formalization can contain mistakes. This site has not independently checked the proofs; the label records the scope of OpenAI's claim.

\subsection{Sources}
\begin{itemize}
\item OpenAI, \emph{A dimension-free logarithmic Sobolev inequality for subgaussian log-concave measures}, 2026-09-23; citation key \texttt{OAI:A-dimension-free-logarithmic-Sobolev-inequality-for-subgaussian-log-concave-measures-September-23-2026}. \href{https://github.com/openai/math/blob/adc7f1241b42e322a6451854ab7e4b4c146bf78a/preprints/A-dimension-free-logarithmic-Sobolev-inequality-for-subgaussian-log-concave-measures-September-23-2026/paper.pdf}{PDF (GitHub)}; \href{https://github.com/openai/math/raw/adc7f1241b42e322a6451854ab7e4b4c146bf78a/preprints/A-dimension-free-logarithmic-Sobolev-inequality-for-subgaussian-log-concave-measures-September-23-2026/paper.pdf}{PDF (direct)}; \href{https://github.com/openai/math/blob/adc7f1241b42e322a6451854ab7e4b4c146bf78a/preprints/A-dimension-free-logarithmic-Sobolev-inequality-for-subgaussian-log-concave-measures-September-23-2026/README.md}{Manuscript page and citation}; \href{https://github.com/openai/math/blob/main/preprints/A-dimension-free-logarithmic-Sobolev-inequality-for-subgaussian-log-concave-measures-September-23-2026/paper.pdf}{Latest version}; Statement source: \href{https://github.com/openai/math/blob/adc7f1241b42e322a6451854ab7e4b4c146bf78a/preprints/A-dimension-free-logarithmic-Sobolev-inequality-for-subgaussian-log-concave-measures-September-23-2026/build/sections/01-introduction.tex\#L26}{thm:main}.
\item \href{https://github.com/openai/math/blob/adc7f1241b42e322a6451854ab7e4b4c146bf78a/CONTENTS.md}{OpenAI catalogue entry}.
\end{itemize}
\end{document}
